% TOP_PROBLEM: {"id": 5300077, "title": "Lift-family criterion for finite kneading data", "category_id": 11, "category": {"id": 11, "name": "dynamical_systems", "display_name": "Dynamical Systems", "description": "Problems about long-term behavior of deterministic systems, Hamiltonian dynamics, and periodic orbits.", "slug": "dynamical-systems", "order_index": 11, "created_at": "2026-07-31T15:26:25.671Z"}, "status": "open"}
% FIRST_POSTED: null
% ENTRY_KIND: statement_only
% Imported from: batch06.tex; UnsolvedMath ID 5300077
\documentclass[11pt]{article}
\usepackage[margin=1in]{geometry}
\usepackage{amsmath,amssymb,mathtools}
\usepackage{fontspec,unicode-math}
\setmainfont{Latin Modern Roman}
\setsansfont{Latin Modern Sans}
\setmathfont{Latin Modern Math}
\usepackage{xurl,hyperref}
\hypersetup{colorlinks=true,linkcolor=blue,urlcolor=blue}
\newcommand{\sref}[2]{\hyperlink{source:#1:#2}{[S#2]}}
\newcommand{\cataloguescope}[1]{\paragraph{Mathematical question.}#1}
\begin{document}
% UnsolvedMath ID 5300077; source catalogue code AMR-052-0077
\setcounter{section}{5300076}
\section{Lift-family criterion for finite kneading data}\label{5300077}
{\small\sffamily UnsolvedMath ID 5300077 · Dynamical Systems}
\subsection{Definitions and mathematical statement}

\paragraph{Shared notation.}$\mathbb N=\{1,2,\ldots\}$, and $\mathbb Z,\mathbb Q,\mathbb R,\mathbb C$ denote the integers, rationals, reals and complex numbers. Any other conventions are specified locally.


Let $I=[0,1]$ and fix the number $n$ of laps (maximal monotonicity intervals) and a boundary map $\mu:\{0,1\}\to\{0,1\}$. Write $\mathcal P(n,\mu)$ for the continuous piecewise monotone interval maps with those data. Starting with $f_0$, lift $f_k$ to the degree-$n$ polynomial $p_k\in\mathcal P(n,\mu)$ whose critical points lie in $I$ and whose ordered critical values equal those of $f_k$. Let $h_k:I\to I$ be the homeomorphism taking ordered critical points of $f_k$ to those of $p_k$ with $p_k\circ h_k=f_k$, and put
\[f_{k+1}=h_k\circ p_k=h_k\circ f_k\circ h_k^{-1}.\]
This is the real Thurston algorithm in the source. A kneading sequence records itineraries of critical values relative to the ordered laps and their boundaries. A map is postcritically finite when its critical points have finite forward orbits. For a different lifting family, replace the polynomial lifts by members of that family having the same ordered critical data and the analogous homeomorphisms.
\paragraph{Statement recorded in the supplied source.}
Find a general property of a lifting family that guarantees convergence of the real Thurston algorithm for every periodic or preperiodic kneading sequence. \sref{5300077}{2}
\subsection{Short English statement}
Find a general criterion on lift families that guarantees convergence for all periodic and preperiodic kneading data.
\subsection{Sources}
\begin{description}
\item[\hypertarget{source:5300077:1}{S1}] ULAM.AI and the UnsolvedMath Contributors, UnsolvedMath: A Curated Collection of Open Mathematics Problems, record 5300077 (AMR-052-0077). CC BY 4.0. \url{https://huggingface.co/datasets/ulamai/UnsolvedMath}.
\item[\hypertarget{source:5300077:2}{S2}] Ben Bielefeld (editor) and conference contributors, \emph{Conformal Dynamics Problem List}, arXiv:math/9201271v1 (1990), \S3, Question 3, printed p. 9. \url{https://arxiv.org/abs/math/9201271v1}.
\end{description}
\label{end:5300077}
\end{document}
